\section{Unconditional finite even-symplectic normal-core exclusions}
\label{sec:higher-even-symplectic-finite}

\noindent\textbf{Status.} Complete authored proof and fresh complete finite
certificates; root whole-proof/code review and fresh $756$-row,
$271785$-entry rerun pass. The complete mathematical argument and all three character tables,
including the newly added third table, have independent whole-proof and
arithmetic review by /root/g34\_eightfolds. Its separate reconstruction
checks all $756$ rows and $271785$ entries; see the retained
\texttt{HIGHER\_EVEN\_SYMPLECTIC\_FINITE\_PEER\_AUDIT.md}. This section does not use the unresolved Gluck--Magaard (1998)
proper-parabolic estimate appearing in the separate all-rank draft.

\begin{theorem}[D-HES-F: three complete core families]
In the signature-$(3,3)$ sixfold-source problem, suppose the faithful image
$J$ of an irreducible inducing representation has an irreducible normal
subgroup $E$ which is a product of simple groups chosen from
\[
          \operatorname{Sp}_8(2),\quad \operatorname{Sp}_{10}(2),
                         \quad \operatorname{Sp}_{12}(2).
\]
Then the actual period image has dimension less than nine. The conclusion
allows arbitrary normalizing scalars, every compatible permutation of
simple factors, arbitrary multiplicities of the three types, and inducing
character fields larger than the original imaginary quadratic field.
No determinant or reflection assumption is needed.
\end{theorem}

\begin{proof}
First we establish a complete finite representation-theoretic bound. The
ordinary character tables supplied by CTblLib are
\texttt{S8(2)}, \texttt{S10(2)} and \texttt{S12(2)}. The first two
provenance strings name the ATLAS~\cite{SPOR:Atlas,SPOR:CTblLib}. The third attributes the table to
Christoph K\"ohler, computed using group-theoretic and character-theoretic
methods~\cite{HES:S12Data}. We use that official complete-table input as such; we do not claim
that it is an ATLAS table or that its construction has a verified separate
refereed publication. For every nontrivial
irreducible character, and every nonidentity class, the exact certificate
establishes the following table:
\[
\begin{array}{c|r|r|c|r}
\text{group}&\text{all irreducible rows}&\text{minimum degree}&
 \max_{\chi\ne1,x\ne1}|\chi(x)|/\chi(1)&\text{attaining value}\\ \hline
\operatorname{Sp}_8(2)&81&35&3/5&-21\\
\operatorname{Sp}_{10}(2)&198&155&17/31&-85\\
\operatorname{Sp}_{12}(2)&477&651&11/21&-341
\end{array}
\tag{HESF.1}
\]
The attaining values occur in rows of the displayed minimum degrees at
involutions. Every row is real valued. The exporter verifies full row
orthogonality and reality in exact cyclotomic arithmetic. Its retained
input includes every row, every class, class sizes, orders, power maps,
Frobenius--Schur indicators, and the original provenance strings.
The independent arithmetic routine uses only this retained input. For an
entry $\sum_jc_j\zeta_N^j$, it bounds the absolute value by
$\sum_j|c_j|$, comparing rational numbers only. All $271785$ nonidentity
entries of nontrivial rows satisfy the displayed bounds. At each maximum
an actual rational entry attains the bound, proving equality, rather than
merely an upper estimate. The routine also verifies group-order formulas,
class-size and degree-square sums, all degrees, and reality by reduction
modulo the relevant cyclotomic polynomial. Thus the common bound $3/5$
is proved for all representations and all nonidentity elements of these
three groups; no sample or extrapolation is involved.

If $A$ is a finite-order matrix of degree $a$ and an eigenvalue has
multiplicity $m$, then $|\operatorname{tr}A|\ge 2m-a$ by the triangle
inequality. It follows from (HESF.1) that every nonidentity core operator
has projective support
\[
          a-\max_\lambda\dim\ker(A-\lambda I)\ge a/5.
\tag{HESF.2}
\]
The core groups are simple, centerless, have trivial Schur multiplier, and
have no outer automorphisms. The precise all-$n\ge4$ classification facts
for $\operatorname{Sp}_{2n}(2)$ are recorded in Gluck--Magaard,
\emph{Absolute Fourth Moments}, author p.~25~\cite{HES:GM2006}. The complete exceptional
multiplier list in \cite[Chapter~5, Table~1, numbered p.~35]{BI:HHM}
contains only the ranks two and three in this even symplectic family.
These known classification facts do not require the fourth-moment
hypothesis used elsewhere in that paper.

Write $V=\bigotimes_iV_i$ for the restriction of the inducing space to
$E$, with $a_i=\dim V_i$ and $e=\dim V$. Every factor module is nontrivial;
a trivial factor could not belong to the faithful image. A normalizer
operator permutes isomorphic simple factors. After that permutation is
removed, each induced automorphism is inner, because the outer groups
are trivial. Schur's lemma therefore makes the operator a scalar times a
tensor of core operators. No additional projective cover is omitted,
since the multipliers of these groups are trivial.

For factor-preserving tensors, eigenvalue multiplicity distributions
convolve on roots of unity. The largest normalized multiplicity of the
convolution is no greater than the largest normalized multiplicity of any
one factor. For a nontrivial factor-permutation cycle of length $l$, the
factor dimensions are equal, say $a$. Conjugating the tensor coordinates
puts all twists on one coordinate; diagonalize their product. On the
resulting tensor basis the operator is a weighted cyclic permutation.
There are $a$ constant words among $a^l$ basis words. Every other cyclic
orbit has length at least two, and each weighted cycle contributes at
most one copy of a specified eigenvalue. Consequently its largest
multiplicity is at most
\[
           a+(a^l-a)/2\le 3a^l/4\qquad(a\ge2,\ l\ge2).
\tag{HESF.3}
\]
Eigenvalue coincidences do not alter the orbit count. Tensoring with the
other cycles cannot increase this normalized bound. Combining (HESF.2)
and (HESF.3), every nonscalar operator in $J$ has projective support at
least $e/5$. A nonidentity scalar has actual moving rank $e$. Thus every
nonidentity inertia in $J$ has actual moving rank at least $e/5$.

The reality certified for the core characters makes each $V_i$ self-dual.
Therefore $V$ has a nondegenerate $E$-invariant bilinear form, whose line
is unique by irreducibility. The full normalizer preserves this line and
acts on it by a multiplier character $\mu:J\to\mathbf C^\times$.
This character is nontrivial here: if it were trivial, every finite-order
operator in $J$ would have reciprocal eigenvalue pairs, and hence real
trace. Its character field would be real. But the inducing character
field contains the original imaginary quadratic field, by rational
induction and Frobenius reciprocity as in D-BI-01. Only that inclusion,
not equality of the two fields, is being used.

We apply the actual intermediate-cover factorization of D-BI-01 to this
inducing representation. If $h$ is the intermediate base genus and
$g_1,\ldots,g_r$ its nonidentity effective inertias, the six-dimensional
complex multiplicity variation has Euler formula
\[
               6=2(h-1)e+\sum_{j=1}^r\operatorname{rank}(g_j-I).
\tag{HESF.4}
\]
Here $e\ge35$. Hence $h\ge2$ is impossible. If $h=1$, the sum is six,
so $r\le6$, and the moduli dimension of the pointed intermediate curve
is at most six (one if $r=0$). The finite-cover factorization bounds the
actual period image by that dimension. A nine-dimensional image would
therefore require $h=0$, $r\ge12$, and
\[
                   \sum_j\operatorname{rank}(g_j-I)=2e+6.
\tag{HESF.5}
\]
This is a necessary condition from an actual variation factorization;
no sufficiency is inferred from a Hurwitz-space dimension.

The inertia tuple generates $J$ and has product one. Since $\mu$ is
nontrivial, at least two entries have multiplier different from one.
For such an entry its $1$-eigenspace is totally isotropic for the
nondegenerate invariant bilinear form, hence has dimension at most
$e/2$. These two entries cost at least $e$ in (HESF.5). The remaining
at least ten nonidentity entries each cost at least $e/5$. Therefore
\[
                    2e+6\ge e+10e/5=3e,
                    \qquad e\le6,
\tag{HESF.6}
\]
contrary to $e\ge35$. This proves the theorem with the full stated
normalizer, tensor, scalar, and larger-field scope.
\end{proof}

\paragraph{Retained certificate and scope.}
The complete frozen input is \texttt{sp2n2\_character\_inputs.json};
its SHA-256 is \texttt{5d3539d84b4988ac34c2703cbc81369b2840e50e3d74ea3afbb2db31f974e615}.
The exporter is \texttt{export\_sp2n2.g}, and the exact verifier and output
are \texttt{check\_sp2n2.py} and \texttt{check\_sp2n2.json}.
The fresh run checks $756$ rows and $271785$ entries. The original two-table
input is retained separately as \texttt{sp2n2\_character\_inputs\_v1.json}
to identify precisely the input used by the earlier peer computation.
The earlier all-rank draft retains its separate (GM) primary-source gate.
A new elementary replacement is under review; none of its conclusions
is used in this finite theorem.
